% SOURCE_STATEMENT: "Within mathematical models of cosmological singularity formation, what mechanism or initial-data principle explains the emergence of a homogeneous, isotropic Big-Bang singularity with vanishing Weyl curvature from more general initial conditions?"
% Preparation and final self-audit: CONTRIBUTING.md prompts C and D; source/audit artifacts retained in the local physics_statements_2026_10_11 batch.
% !TeX program = lualatex
\documentclass[11pt]{article}
\usepackage[margin=25mm]{geometry}
\usepackage{fontspec}
\setmainfont{Latin Modern Roman}
\usepackage{amsmath,amssymb,mathtools}
\usepackage{unicode-math}
\setmathfont{Latin Modern Math}
\usepackage{xurl,hyperref}
\hypersetup{colorlinks=true,urlcolor=blue}
\setcounter{secnumdepth}{0}
\setlength{\parindent}{0pt}
\setlength{\parskip}{5pt}
\setlength{\emergencystretch}{3em}
\newcommand{\sref}[2]{\hyperlink{source:#1:#2}{[S#2]}}
\newcommand{\cataloguescope}[1]{\paragraph{Recorded target.}#1}
\newcommand{\R}{\mathbb R}
\newcommand{\C}{\mathbb C}
\newcommand{\N}{\mathbb N}
\newcommand{\Z}{\mathbb Z}
\newcommand{\Q}{\mathbb Q}
\newcommand{\F}{\mathbb F}
\newcommand{\eps}{\varepsilon}
\begin{document}
\section{Emergence of Isotropic Big-Bang Singularities}\label{statement}

\subsection{Definitions and mathematical statement}

\paragraph{Shared notation.}$\mathbb N=\{1,2,\ldots\}$; the standard number fields have their usual meanings. All additional conventions are stated locally.


\paragraph{Definitions and conventions.}
A Big-Bang singularity is a past boundary where the classical spacetime description ceases to be regular, often accompanied by curvature blowup or geodesic incompleteness. Homogeneity means that spatial slices have transitive symmetry; isotropy means rotational symmetry about spatial points. The Weyl tensor is the trace-free part of spacetime curvature and vanishes in a conformally flat homogeneous isotropic cosmological model. A mechanism or initial-data principle must explain why this special singular behavior arises from a wider class of data.

\paragraph{Target retained from the canonical source catalogue.}
\paragraph{Catalogue wording retained for traceability.}
Within mathematical models of cosmological singularity formation, what mechanism or initial-data principle explains the emergence of a homogeneous, isotropic Big-Bang singularity with vanishing Weyl curvature from more general initial conditions?

\paragraph{Source scope and qualification.}
This is a cosmological mechanism question rather than one fully specified theorem. The matter equations, admissible data, singularity limit and meaning of emergence are not fixed by the extraction; a uniqueness account of the universe\textquotesingle s origin is not inferred.
\subsection{Short English statement}
Explain mathematically why general cosmological data could produce a homogeneous isotropic initial singularity with vanishing Weyl curvature.
\subsection{Sources}
\begin{itemize}
\item[\hypertarget{source:statement:1}{S1}] ULAM.AI and the UnsolvedMath Contributors, UnsolvedMath, catalogue entry. Catalogue curation under CC BY 4.0; definitions and English exposition added locally.
\url{https://huggingface.co/datasets/ulamai/UnsolvedMath}.
\item[\hypertarget{source:statement:2}{S2}] Mathematical Aspects of General Relativity (2021), Oberwolfach Reports, DOI 10.4171/owr/2021/40, pp. 2157--2267. Primary report accessed 2026-10-11; relevant PDF pages 26, 27.
\url{https://doi.org/10.4171/owr/2021/40}.
\item[\hypertarget{source:statement:3}{S3}] Borde--Guth--Vilenkin, inflationary incompleteness (catalogue-listed context; not independently consulted in this preparation).
\url{https://arxiv.org/abs/gr-qc/0110012}.
\end{itemize}
\label{end:statement}
\end{document}
